%=================================================
% CHAPITRE 1 : Cadre général du projet
%=================================================

\chapter{Cadre général du projet}

\section{Introduction}

Le point de départ de ce projet est un constat assez simple, mais très concret : dans beaucoup d'administrations scolaires ou universitaires, la gestion des ressources humaines ressemble encore à un puzzle de fichiers Excel et de documents Word copiés-collés. On passe un temps fou à ressaisir les mêmes noms, les mêmes grades et les mêmes dates dans des attestations de travail ou des contrats. 

L'idée de SIRH-Doc est née de cette frustration. L'objectif n'est pas seulement de stocker des noms dans une base de données, mais de créer une plateforme intelligente capable de transformer ces données en documents officiels, tout en gardant une trace impeccable de chaque décision administrative. Ce premier chapitre pose le décor : pourquoi ce projet, pour qui, et avec quels outils j'ai décidé de le construire.

\section{Présentation de l'organisme d'accueil}

Le projet s'inscrit dans le cadre d'un établissement d'enseignement supérieur. C'est un environnement où la gestion du personnel est particulièrement complexe : il faut jongler entre les enseignants permanents, les agents administratifs et une multitude de vacataires dont les contrats changent à chaque session.

Chaque jour, le service RH doit produire des dizaines de pièces : certificats d'enseignement, états de services, attestations de présence... La moindre erreur de frappe sur un grade ou une date peut invalider un document officiel. L'établissement disposait déjà de données informatisées, mais il manquait ce \og pont \fg{} automatisé entre la donnée brute et le document final prêt à imprimer.

\section{Présentation du projet}

\subsection{Contexte et Problématique}

Aujourd'hui, quand on veut créer une attestation pour un enseignant, on cherche souvent le dernier fichier Word qui a servi de modèle, on fait un \og Enregistrer sous \fg{}, et on modifie manuellement les informations. Cette méthode artisanale pose trois problèmes majeurs :
\begin{itemize}
    \item \textbf{Le risque d'erreur} : On oublie souvent de changer une date ou un titre caché en bas de page.
    \item \textbf{La perte d'historique} : Il est très difficile de savoir exactement ce qui a été délivré à qui il y a deux ans.
    \item \textbf{L'hétérogénéité} : Selon la personne qui rédige, le logo n'est pas à la même place ou la police change.
\end{itemize}

Ma problématique était donc la suivante : \textit{Comment centraliser la gestion du personnel et automatiser la production documentaire pour garantir à la fois un gain de temps et une fiabilité totale ?}

\subsection{La solution : SIRH-Doc}

La solution que j'ai conçue est une plateforme web intégrée. Elle ne se contente pas d'être un éditeur de texte ; c'est un véritable hub où l'administrateur définit des \og familles de documents \fg{}, configure des modèles dynamiques avec des variables (comme \texttt{\{\{nom\}\}} ou \texttt{\{\{grade\}\}}), et laisse le système fusionner tout cela avec les données réelles du personnel.

\section{Méthodologie de travail}

Même si j'ai mené ce développement seul, j'ai voulu suivre une démarche structurée en m'inspirant de la méthode **Agile Scrum**. J'ai découpé le projet en quatre sprints de trois à quatre semaines chacun. 

Cette approche m'a permis de ne pas me laisser submerger par la complexité de l'éditeur documentaire dès le début. J'ai d'abord sécurisé les fondations (Sprint 1), puis construit le moteur de données (Sprint 2), avant de m'attaquer au studio d'édition (Sprint 3) et à la production finale (Sprint 4).

\section{Environnement technique : Mes choix}

\subsection{Backend : .NET 8 et Dapper}
Pour le serveur, j'ai choisi **ASP.NET Core (Web API)**. C'est un framework que j'apprécie pour sa rigueur et sa performance. Pour l'accès aux données, j'ai préféré **Dapper** à Entity Framework. Pourquoi ? Parce que SIRH-Doc doit pouvoir exécuter des requêtes SQL très spécifiques et dynamiques selon les besoins des établissements, et Dapper me donne ce contrôle total sur le SQL tout en restant extrêmement rapide.

\subsection{Frontend : Angular 17}
Côté interface, **Angular** s'est imposé naturellement. Sa structure modulaire est parfaite pour une application d'administration lourde. J'ai utilisé la version 17 pour profiter des dernières optimisations de performance (comme le nouveau système de contrôle de flux).

\subsection{L'édition : Tiptap}
Pour la partie \og Studio \fg{}, j'ai intégré **Tiptap**. Contrairement aux éditeurs classiques qui sont parfois trop rigides, Tiptap est \og headless \fg{}. Cela m'a permis de construire une interface de création de modèles totalement personnalisée, tout en produisant un code HTML propre pour la génération des PDF.

\section{Conclusion}

Ce chapitre a permis de définir le \og pourquoi \fg{} du projet. SIRH-Doc n'est pas un simple exercice technique, c'est une réponse à un besoin de modernisation administrative. Maintenant que le cadre est posé, le chapitre suivant va détailler plus finement les fonctionnalités attendues et la manière dont j'ai modélisé le système.
